\chapter{A Safety Case for Superintelligence Alignment}
\label{ch:safety-case}

\begin{chapterthesis}
\textbf{[STUB]} One paragraph stating the chapter's core claim.
\end{chapterthesis}

\begin{refsection}

\epigraph{%
The crucial variables are discovered by accident.%
}{--- John Gall, \emph{Systemantics} (1975)}

\section{Why This Matters}

\textbf{[STUB]} Explain the decision relevance.

\section{Plain-Language Model}

\textbf{[STUB]} Explain the idea without equations.

\section{Formal Model}

\textbf{[STUB]} Introduce variables, assumptions, and equations.

\section{Worked Example}

\textbf{[STUB]} Give one concrete example.

\section{Counterexample or Failure Mode}

\textbf{[STUB]} Give at least one case where the model fails, overreaches, or needs refinement.

\section{What Would Change This View}

\begin{itemize}
  \item A complete, rigorous safety case---every claim discharged, every assumption checked---is built for a system that then causes catastrophe, because a safety case can only certify the absence of \emph{imagined} failure modes and a superintelligence supplies an unimagined one.
  \item The case rests on metrics that are observable but not adversarially verifiable, so a passing case and a deceptive system are compatible (Chapter~\ref{ch:verifiability-ontology}). What would strengthen it: a safety-case format whose failure the system is forced to reveal \emph{before} the loss, not one that explains the loss afterward.
\end{itemize}

\section{Summary}

\begin{itemize}
  \item \textbf{[STUB]} Summary point 1.
  \item \textbf{[STUB]} Summary point 2.
  \item \textbf{[STUB]} Summary point 3.
\end{itemize}

\section*{Chapter References}

This chapter builds on safety-case and assurance-argument methodology \autocite{kelly1998safety,bloomfield2012safety}; frontier-safety commitments \autocite{seoul2024commitments}; and the open-risk synthesis of the International AI Safety Report \autocite{iaisr2025}.

\printbibliography[heading=subbibliography,title={Chapter References}]

\end{refsection}

